\chapter{Abstract}
Dieses Projekt ist ganz toll, ich schwör.\\

It is remarkable that all ordinary visible matter is composed of only three stable fundamental particles: the electron, the up quark, and the down quark. While other elementary particles, such as photons, gluons, and the weak gauge bosons, play essential roles in mediating the fundamental interactions, they do not form stable matter themselves. Despite the remarkable success of the Standard Model, it is widely recognized as an incomplete theory, as it cannot explain several fundamental phenomena, such as the nature of dark matter, the origin and ordering of neutrino masses, or the matter-antimatter asymmetry observed in the Universe.
Among the many open questions in particle physics is the precise understanding of the proton's internal structure. The so-called proton radius puzzle emerged from conflicting measurements of the proton's charge radius obtained using different experimental approaches. Although recent experimental results have substantially reduced this discrepancy, an independent high-precision determination of the proton radius remains of great importance. The AMBER experiment at CERN aims to contribute to this effort by measuring the proton charge radius with unprecedented precision.